\clearpage
\begingroup
\fancyhead[L]{\small INHALT.}
\renewcommand{\headrulewidth}{0pt}
\begingroup
\newcommand{\ContentsEntry}[4]{\par\begingroup\leftskip=#1\relax\rightskip=0pt plus 1fil\parfillskip=0pt\hangindent=1.5em\hangafter=1\noindent#2\nobreak\hspace{.4em}\mbox{#4}\nobreak\hspace{.5em}\hfill\mbox{#3}\par\endgroup}
\newcommand{\ContentsPoint}{\hbox to .38em{\hfil.\hfil}}
\newcommand{\ContentsText}[2]{\par\begingroup\leftskip=#1\relax\hangindent=1.5em\hangafter=1\noindent#2\par\endgroup}
\newcommand{\ContentsHeading}[1]{\needspace{4\baselineskip}\begin{center}#1\end{center}}
\newcommand{\ContentsRule}{\begin{center}\rule{.14\linewidth}{.4pt}\end{center}}
\newcommand{\ContentsJupiter}{\text{\mdseries\AstroSymbols ♃}}
\newcommand{\ContentsPallas}{\text{\mdseries\AstroSymbols ⚴}}
\SourceLeaf{650}
\needspace{10\baselineskip}
\pdfbookmark[1]{Inhalt}{vii-historical-inhalt}
\ContentsHeading{{\Large INHALT.}}
\ContentsHeading{{\large GAUSS WERKE BAND VII.}\\THEORETISCHE ASTRONOMIE.}
\ContentsRule
\ContentsHeading{THEORIA MOTUS CORPORUM COELESTIUM IN SECTIONIBUS\\CONICIS SOLEM AMBIENTIUM.}
\ContentsEntry{1.5em}{Praefatio}{Seite 3}{\ContentsPoint\ContentsPoint\ContentsPoint\ContentsPoint\ContentsPoint\ContentsPoint\ContentsPoint\ContentsPoint\ContentsPoint\ContentsPoint\ContentsPoint\ContentsPoint\ContentsPoint\ContentsPoint\ContentsPoint\ContentsPoint\ContentsPoint\ContentsPoint\ContentsPoint\ContentsPoint\ContentsPoint\ContentsPoint\ContentsPoint\ContentsPoint\ContentsPoint\ContentsPoint\ContentsPoint\ContentsPoint\ContentsPoint\ContentsPoint}
\ContentsText{1.5em}{Liber primus. Relationes generales inter quantitates, per quas corporum coelestium motus circa solem definiuntur.}
\ContentsEntry{3em}{Sectio I. Relationes ad locum simplicem in orbita spectantes}{--- 13}{\ContentsPoint\ContentsPoint\ContentsPoint\ContentsPoint\ContentsPoint\ContentsPoint\ContentsPoint\ContentsPoint\ContentsPoint}
\ContentsEntry{3em}{Sectio II. Relationes ad locum simplicem in spatio spectantes}{--- 59}{\ContentsPoint\ContentsPoint\ContentsPoint\ContentsPoint\ContentsPoint\ContentsPoint\ContentsPoint\ContentsPoint\ContentsPoint}
\ContentsEntry{3em}{Sectio III. Relationes inter locos plures in orbita}{--- 100}{\ContentsPoint\ContentsPoint\ContentsPoint\ContentsPoint\ContentsPoint\ContentsPoint\ContentsPoint\ContentsPoint\ContentsPoint\ContentsPoint\ContentsPoint\ContentsPoint\ContentsPoint}
\ContentsEntry{3em}{Sectio IV. Relationes inter locos plures in spatio}{--- 148}{\ContentsPoint\ContentsPoint\ContentsPoint\ContentsPoint\ContentsPoint\ContentsPoint\ContentsPoint\ContentsPoint\ContentsPoint\ContentsPoint\ContentsPoint\ContentsPoint\ContentsPoint}
\ContentsText{1.5em}{Liber secundus. Investigatio orbitarum corporum coelestium ex observationibus geocentricis.}
\ContentsEntry{3em}{Sectio I. Determinatio orbitae e tribus observationibus completis}{--- 155}{\ContentsPoint\ContentsPoint\ContentsPoint\ContentsPoint\ContentsPoint\ContentsPoint\ContentsPoint\ContentsPoint}
\ContentsEntry{3em}{Sectio II. Determinatio orbitae e quatuor observationibus, quarum duae tantum completae sunt}{--- 222}{\ContentsPoint\ContentsPoint\ContentsPoint\ContentsPoint\ContentsPoint\ContentsPoint\ContentsPoint\ContentsPoint\ContentsPoint\ContentsPoint\ContentsPoint\ContentsPoint\ContentsPoint\ContentsPoint\ContentsPoint\ContentsPoint\ContentsPoint\ContentsPoint\ContentsPoint\ContentsPoint\ContentsPoint\ContentsPoint\ContentsPoint\ContentsPoint\ContentsPoint\ContentsPoint}
\ContentsEntry{3em}{Sectio III. Determinatio orbitae observationibus quotcunque quam proxime satisfacientis}{--- 236}{}
\ContentsEntry{3em}{Sectio IV. De determinatione orbitarum, habita ratione perturbationum}{--- 258}{\ContentsPoint\ContentsPoint\ContentsPoint\ContentsPoint\ContentsPoint\ContentsPoint}
\ContentsEntry{1.5em}{Tabulae}{--- 263}{\ContentsPoint\ContentsPoint\ContentsPoint\ContentsPoint\ContentsPoint\ContentsPoint\ContentsPoint\ContentsPoint\ContentsPoint\ContentsPoint\ContentsPoint\ContentsPoint\ContentsPoint\ContentsPoint\ContentsPoint\ContentsPoint\ContentsPoint\ContentsPoint\ContentsPoint\ContentsPoint\ContentsPoint\ContentsPoint\ContentsPoint\ContentsPoint\ContentsPoint\ContentsPoint\ContentsPoint\ContentsPoint\ContentsPoint\ContentsPoint}
\ContentsEntry{0pt}{Bemerkungen}{--- 281}{\ContentsPoint\ContentsPoint\ContentsPoint\ContentsPoint\ContentsPoint\ContentsPoint\ContentsPoint\ContentsPoint\ContentsPoint\ContentsPoint\ContentsPoint\ContentsPoint\ContentsPoint\ContentsPoint\ContentsPoint\ContentsPoint\ContentsPoint\ContentsPoint\ContentsPoint\ContentsPoint\ContentsPoint\ContentsPoint\ContentsPoint\ContentsPoint\ContentsPoint\ContentsPoint\ContentsPoint\ContentsPoint\ContentsPoint\ContentsPoint}
\ContentsRule
\ContentsHeading{ZUSÄTZE ZUR THEORIA MOTUS CORPORUM COELESTIUM\\UND ANDERE NACHTRÄGE ZUR ELLIPTISCHEN BEWEGUNG.}
\ContentsText{0pt}{I. Zusätze zur Theoria motus.}
\ContentsText{1.5em}{\textit{Nachlass.}}
\ContentsEntry{1.5em}{Zu Art. 7. Genäherte Formeln für die grösste Mittelpunktsgleichung}{--- 291}{\ContentsPoint\ContentsPoint\ContentsPoint\ContentsPoint\ContentsPoint\ContentsPoint\ContentsPoint\ContentsPoint}
\ContentsEntry{1.5em}{Zu Art. 53. Projection einer Planetenbahn auf die Ekliptik}{--- 291}{\ContentsPoint\ContentsPoint\ContentsPoint\ContentsPoint\ContentsPoint\ContentsPoint\ContentsPoint\ContentsPoint\ContentsPoint\ContentsPoint\ContentsPoint}

\SourceLeaf{651}
\needspace{4\baselineskip}
\noindent\hfill INHALT.\hfill 645\par
\ContentsEntry{1.5em}{Zu Art. 70. Zur Berechnung der Parallaxe}{Seite 295}{\ContentsPoint\ContentsPoint\ContentsPoint\ContentsPoint\ContentsPoint\ContentsPoint\ContentsPoint\ContentsPoint\ContentsPoint\ContentsPoint\ContentsPoint\ContentsPoint\ContentsPoint\ContentsPoint\ContentsPoint\ContentsPoint\ContentsPoint}
\ContentsEntry{1.5em}{Zu Art. 76--77. Allgemeine Differentialformeln für die geocentrischen Örter der Planeten}{--- 296}{\ContentsPoint}
\ContentsEntry{1.5em}{Zu Art. 78. Directe Auflösung der Gleichungen \(\sin(P-A)=k\sin(Q-B),\ \sin(P-A')=k'\sin(Q-B')\)}{--- 297}{\ContentsPoint\ContentsPoint\ContentsPoint\ContentsPoint\ContentsPoint\ContentsPoint\ContentsPoint\ContentsPoint\ContentsPoint\ContentsPoint\ContentsPoint\ContentsPoint\ContentsPoint\ContentsPoint\ContentsPoint\ContentsPoint\ContentsPoint\ContentsPoint\ContentsPoint\ContentsPoint\ContentsPoint\ContentsPoint\ContentsPoint\ContentsPoint\ContentsPoint}
\ContentsEntry{1.5em}{Zu Art. 88--90. Näherungsformel, um aus \(\sin\frac14\varphi\) den Logarithmen von \(\dfrac{\varphi-\sin\varphi}{\frac43\sin\frac12\varphi^3}\) zu finden; Vorschriften, um den Logar. Sinus eines kleinen Bogens zu finden}{--- 299}{\ContentsPoint\ContentsPoint\ContentsPoint\ContentsPoint\ContentsPoint\ContentsPoint\ContentsPoint}
\ContentsText{1.5em}{\textit{Veröffentlichungen.}}
\ContentsEntry{1.5em}{Zu Art. 90 und 100. Zur Auflösung der Aufgabe, aus zweien Radiis vectoribus und dem eingeschlossenen Winkel die elliptischen oder hyperbolischen Elemente zu bestimmen}{--- 300}{\ContentsPoint}
\ContentsEntry{1.5em}{Zu Art. 114 und 177. Druckfehler in Dr. \textsc{Gauss}’ Theoria motus etc. vom Senator Bar. \textsc{Oriani}; über das Integral \(\int\limits_0^\infty e^{-tt}\mathrm{d}t=\frac12\BareRoot{\pi}\)}{--- 301}{\ContentsPoint\ContentsPoint\ContentsPoint\ContentsPoint\ContentsPoint\ContentsPoint\ContentsPoint\ContentsPoint\ContentsPoint\ContentsPoint\ContentsPoint\ContentsPoint\ContentsPoint\kern.38em\ContentsPoint}
\ContentsText{1.5em}{\textit{Nachlass.}}
\ContentsEntry{1.5em}{Zu Art. 138. Criterium des Falls, wo der die geocentrische Bewegung tangirende grösste Kreis durch den Sonnenort geht}{--- 303}{\ContentsPoint\ContentsPoint\ContentsPoint\ContentsPoint\ContentsPoint\ContentsPoint\ContentsPoint\ContentsPoint\ContentsPoint\ContentsPoint\ContentsPoint\ContentsPoint\ContentsPoint\ContentsPoint\ContentsPoint\ContentsPoint\ContentsPoint\ContentsPoint\ContentsPoint\ContentsPoint}
\ContentsEntry{1.5em}{Zu Art. 140. Ableitung einiger Gleichungen dieses Artikels}{--- 304}{\ContentsPoint\ContentsPoint\ContentsPoint\ContentsPoint\ContentsPoint\ContentsPoint\ContentsPoint\ContentsPoint\ContentsPoint\ContentsPoint\ContentsPoint}
\ContentsText{1.5em}{\textit{Veröffentlichung.}}
\ContentsEntry{1.5em}{Zu Art. 182--184. Berichtigungen}{--- 307}{\ContentsPoint\ContentsPoint\ContentsPoint\ContentsPoint\ContentsPoint\ContentsPoint\ContentsPoint\ContentsPoint\ContentsPoint\ContentsPoint\ContentsPoint\ContentsPoint\ContentsPoint\ContentsPoint\ContentsPoint\ContentsPoint\ContentsPoint\ContentsPoint\ContentsPoint\ContentsPoint}
\ContentsText{1.5em}{\textit{Nachlass.}}
\ContentsEntry{1.5em}{Zu Art. 183. Bestimmung der wahrscheinlichen Fehler der Resultate der Methode der kleinsten Quadrate}{--- 307}{\ContentsPoint\ContentsPoint\ContentsPoint\ContentsPoint\ContentsPoint\ContentsPoint\ContentsPoint\ContentsPoint\ContentsPoint\ContentsPoint\ContentsPoint\ContentsPoint\ContentsPoint\ContentsPoint\ContentsPoint\ContentsPoint\ContentsPoint\ContentsPoint\ContentsPoint\ContentsPoint\ContentsPoint\ContentsPoint\ContentsPoint\ContentsPoint\ContentsPoint}
\ContentsEntry{1.5em}{Zusammenhang des Initialzustandes mit den Elementen einer Planetenbahn}{--- 309}{\ContentsPoint\ContentsPoint\ContentsPoint\ContentsPoint\ContentsPoint\ContentsPoint}
\ContentsEntry{1.5em}{Über die Verbesserung der Elemente eines Planeten durch vier beobachtete Oppositionen}{--- 310}{\ContentsPoint}
\ContentsEntry{1.5em}{Differentialformeln bei Berechnung der Oppositionen}{--- 311}{\ContentsPoint\ContentsPoint\ContentsPoint\ContentsPoint\ContentsPoint\ContentsPoint\ContentsPoint\ContentsPoint\ContentsPoint\ContentsPoint\ContentsPoint\ContentsPoint\ContentsPoint\ContentsPoint}
\ContentsText{0pt}{II. Über die Zodiaken der Himmelskörper.}
\ContentsText{1.5em}{\textit{Veröffentlichung.}}
\ContentsEntry{1.5em}{Schreiben des Herrn Geheimen Hofraths \textsc{Gauss} an den Herausgeber der Astronomischen Nachrichten}{--- 313}{\ContentsPoint\ContentsPoint\ContentsPoint\ContentsPoint\ContentsPoint\ContentsPoint\ContentsPoint\ContentsPoint\ContentsPoint\ContentsPoint\ContentsPoint\ContentsPoint\ContentsPoint\ContentsPoint\ContentsPoint\ContentsPoint\ContentsPoint\ContentsPoint\ContentsPoint\ContentsPoint\ContentsPoint\ContentsPoint\ContentsPoint\ContentsPoint\ContentsPoint\ContentsPoint\ContentsPoint}
\ContentsText{1.5em}{\textit{Nachlass.}}
\ContentsEntry{1.5em}{Zodiakus auf der Himmelskugel}{--- 316}{\ContentsPoint\ContentsPoint\ContentsPoint\ContentsPoint\ContentsPoint\ContentsPoint\ContentsPoint\ContentsPoint\ContentsPoint\ContentsPoint\ContentsPoint\ContentsPoint\ContentsPoint\ContentsPoint\ContentsPoint\ContentsPoint\ContentsPoint\ContentsPoint\ContentsPoint\ContentsPoint\ContentsPoint}
\ContentsEntry{1.5em}{Zodiakus}{--- 319}{\ContentsPoint\ContentsPoint\ContentsPoint\ContentsPoint\ContentsPoint\ContentsPoint\ContentsPoint\ContentsPoint\ContentsPoint\ContentsPoint\ContentsPoint\ContentsPoint\ContentsPoint\ContentsPoint\ContentsPoint\ContentsPoint\ContentsPoint\ContentsPoint\ContentsPoint\ContentsPoint\ContentsPoint\ContentsPoint\ContentsPoint\ContentsPoint\ContentsPoint\ContentsPoint\ContentsPoint\ContentsPoint\ContentsPoint\ContentsPoint}
\ContentsEntry{0pt}{Bemerkungen}{--- 322}{\ContentsPoint\ContentsPoint\ContentsPoint\ContentsPoint\ContentsPoint\ContentsPoint\ContentsPoint\ContentsPoint\ContentsPoint\ContentsPoint\ContentsPoint\ContentsPoint\ContentsPoint\ContentsPoint\ContentsPoint\ContentsPoint\ContentsPoint\ContentsPoint\ContentsPoint\ContentsPoint\ContentsPoint\ContentsPoint\ContentsPoint\ContentsPoint\ContentsPoint\ContentsPoint\ContentsPoint\ContentsPoint\ContentsPoint\ContentsPoint}
\ContentsRule
\ContentsHeading{ZUR PARABOLISCHEN BEWEGUNG.}
\ContentsText{0pt}{I. Zur Cometenbahnbestimmung.}
\ContentsText{1.5em}{\textit{Nachlass.}}
\ContentsEntry{1.5em}{Parabolische Bewegung. Zeit durch zwei Radios vectores und Chorde}{--- 325}{\ContentsPoint\ContentsPoint\ContentsPoint\ContentsPoint\ContentsPoint\ContentsPoint\ContentsPoint\ContentsPoint}
\ContentsEntry{1.5em}{Zur Berechnung der Chorde aus der Zeit und der Summe der beiden Radien}{--- 327}{\ContentsPoint\ContentsPoint\ContentsPoint\ContentsPoint\ContentsPoint}
\ContentsEntry{1.5em}{Zur parabolischen Bewegung. Bestimmung der Elemente durch zwei Radios vectores und Chorde}{--- 327}{\ContentsPoint\ContentsPoint\ContentsPoint\ContentsPoint\ContentsPoint\ContentsPoint\ContentsPoint\ContentsPoint\ContentsPoint\ContentsPoint\ContentsPoint\ContentsPoint\ContentsPoint\ContentsPoint\ContentsPoint\ContentsPoint\ContentsPoint\ContentsPoint\ContentsPoint\ContentsPoint\ContentsPoint\ContentsPoint\ContentsPoint\ContentsPoint\ContentsPoint\ContentsPoint\ContentsPoint\ContentsPoint\ContentsPoint}
\ContentsEntry{1.5em}{Verhältniss von Dreieck und Sector zwischen zwei Radiis vectoribus in der parabolischen Bewegung}{--- 329}{\ContentsPoint\ContentsPoint\ContentsPoint\ContentsPoint\ContentsPoint\ContentsPoint\ContentsPoint\ContentsPoint\ContentsPoint\ContentsPoint\ContentsPoint\ContentsPoint\ContentsPoint\ContentsPoint\ContentsPoint\ContentsPoint\ContentsPoint\ContentsPoint\ContentsPoint\ContentsPoint\ContentsPoint\ContentsPoint\ContentsPoint\ContentsPoint\ContentsPoint\ContentsPoint\ContentsPoint\ContentsPoint}

\SourceLeaf{652}
\needspace{4\baselineskip}
\noindent 646\hfill INHALT.\hfill\mbox{}\par
\ContentsText{1.5em}{\textit{Briefwechsel.}}
\ContentsEntry{1.5em}{Annäherungsformel für die Zwischenzeit: \textsc{Gauss} an \textsc{Olbers} 1815 Januar 7}{Seite 332}{\ContentsPoint\ContentsPoint\ContentsPoint\ContentsPoint\ContentsPoint\ContentsPoint}
\ContentsEntry{13em}{\textsc{Gauss} an \textsc{Olbers} 1815 Januar 13}{--- 335}{\ContentsPoint\ContentsPoint\ContentsPoint\ContentsPoint\ContentsPoint\ContentsPoint}
\ContentsText{1.5em}{\textit{Nachlass.}}
\ContentsEntry{1.5em}{Allgemeine Theorie der Berechnung der Cometenbahnen, nebst Anwendung auf den zweiten Cometen des Jahrs 1813}{--- 338}{\ContentsPoint\ContentsPoint\ContentsPoint\ContentsPoint\ContentsPoint\ContentsPoint\ContentsPoint\ContentsPoint\ContentsPoint\ContentsPoint\ContentsPoint\ContentsPoint\ContentsPoint\ContentsPoint\ContentsPoint\ContentsPoint\ContentsPoint\ContentsPoint\ContentsPoint\ContentsPoint\ContentsPoint\ContentsPoint\ContentsPoint}
\ContentsText{1.5em}{\textit{Briefwechsel.}}
\ContentsText{1.5em}{Supplement zur Theoria motus corporum coelestium:}
\ContentsEntry{10em}{\textsc{Gauss} an \textsc{Olbers} 1815 Mai 29}{--- 350}{\ContentsPoint\ContentsPoint\ContentsPoint\ContentsPoint\ContentsPoint\ContentsPoint\ContentsPoint\ContentsPoint\ContentsPoint}
\ContentsEntry{10em}{\textsc{Gauss} an \textsc{Olbers} 1815 Juni 24}{--- 350}{\ContentsPoint\ContentsPoint\ContentsPoint\ContentsPoint\ContentsPoint\ContentsPoint\ContentsPoint\ContentsPoint\ContentsPoint}
\ContentsText{0pt}{II. Zur Berechnung der wahren Anomalie in der parabolischen Bewegung.}
\ContentsText{1.5em}{\textit{Nachlass.}}
\ContentsEntry{1.5em}{Tabula nova motus parabolici}{--- 351}{\ContentsPoint\ContentsPoint\ContentsPoint\ContentsPoint\ContentsPoint\ContentsPoint\ContentsPoint\ContentsPoint\ContentsPoint\ContentsPoint\ContentsPoint\ContentsPoint\ContentsPoint\ContentsPoint\ContentsPoint\ContentsPoint\ContentsPoint\ContentsPoint\ContentsPoint\ContentsPoint\ContentsPoint\ContentsPoint}
\ContentsText{1.5em}{\textit{Briefwechsel.}}
\ContentsText{1.5em}{Zur Berechnung der Tafel für die parabolische Bewegung:}
\ContentsEntry{10em}{\textsc{Gauss} an \textsc{Encke} 1815}{--- 368}{\ContentsPoint\ContentsPoint\ContentsPoint\ContentsPoint\ContentsPoint\ContentsPoint\ContentsPoint\ContentsPoint\ContentsPoint\ContentsPoint\ContentsPoint\ContentsPoint}
\ContentsEntry{10em}{\textsc{Gauss} an \textsc{Olbers} 1816 Februar 16}{--- 370}{\ContentsPoint\ContentsPoint\ContentsPoint\ContentsPoint\ContentsPoint\ContentsPoint\ContentsPoint\ContentsPoint}
\ContentsText{1.5em}{\textit{Veröffentlichung.}}
\ContentsEntry{1.5em}{Schreiben des Herrn Hofraths \textsc{Gauss} an den Herausgeber der Astronomischen Nachrichten 1843 April}{--- 371}{\ContentsPoint\ContentsPoint\ContentsPoint\ContentsPoint\ContentsPoint\ContentsPoint\ContentsPoint\ContentsPoint\ContentsPoint\ContentsPoint\ContentsPoint\ContentsPoint\ContentsPoint\ContentsPoint\ContentsPoint\ContentsPoint\ContentsPoint\ContentsPoint\ContentsPoint\ContentsPoint\ContentsPoint\ContentsPoint\ContentsPoint\ContentsPoint\ContentsPoint\ContentsPoint\ContentsPoint}
\ContentsEntry{0pt}{Bemerkungen}{--- 373}{\ContentsPoint\ContentsPoint\ContentsPoint\ContentsPoint\ContentsPoint\ContentsPoint\ContentsPoint\ContentsPoint\ContentsPoint\ContentsPoint\ContentsPoint\ContentsPoint\ContentsPoint\ContentsPoint\ContentsPoint\ContentsPoint\ContentsPoint\ContentsPoint\ContentsPoint\ContentsPoint\ContentsPoint\ContentsPoint\ContentsPoint\ContentsPoint\ContentsPoint\ContentsPoint\ContentsPoint\ContentsPoint\ContentsPoint\ContentsPoint}
\ContentsRule
\ContentsHeading{STÖRUNGEN DER CERES.}
\ContentsText{0pt}{I. Erste Methode.}
\ContentsText{1.5em}{\textit{Briefwechsel:}}
\ContentsEntry{1.5em}{\textsc{Gauss} an \textsc{Olbers} 1802 October 12}{--- 377}{\ContentsPoint\ContentsPoint\ContentsPoint\ContentsPoint\ContentsPoint\ContentsPoint\ContentsPoint\ContentsPoint\ContentsPoint\ContentsPoint\ContentsPoint\ContentsPoint\ContentsPoint\ContentsPoint\ContentsPoint\ContentsPoint\ContentsPoint\ContentsPoint\ContentsPoint\ContentsPoint}
\ContentsEntry{1.5em}{\textsc{Gauss} an \textsc{Olbers} 1802 October 26}{--- 378}{\ContentsPoint\ContentsPoint\ContentsPoint\ContentsPoint\ContentsPoint\ContentsPoint\ContentsPoint\ContentsPoint\ContentsPoint\ContentsPoint\ContentsPoint\ContentsPoint\ContentsPoint\ContentsPoint\ContentsPoint\ContentsPoint\ContentsPoint\ContentsPoint\ContentsPoint\ContentsPoint}
\ContentsEntry{1.5em}{\textsc{Gauss} an \textsc{Gerling} 1815 December 28}{--- 378}{\ContentsPoint\ContentsPoint\ContentsPoint\ContentsPoint\ContentsPoint\ContentsPoint\ContentsPoint\ContentsPoint\ContentsPoint\ContentsPoint\ContentsPoint\ContentsPoint\ContentsPoint\ContentsPoint\ContentsPoint\ContentsPoint\ContentsPoint\ContentsPoint\ContentsPoint}
\ContentsText{1.5em}{\textit{Nachlass.}}
\ContentsEntry{1.5em}{Die Störungen der Planeten}{--- 378}{\ContentsPoint\ContentsPoint\ContentsPoint\ContentsPoint\ContentsPoint\ContentsPoint\ContentsPoint\ContentsPoint\ContentsPoint\ContentsPoint\ContentsPoint\ContentsPoint\ContentsPoint\ContentsPoint\ContentsPoint\ContentsPoint\ContentsPoint\ContentsPoint\ContentsPoint\ContentsPoint\ContentsPoint\ContentsPoint\ContentsPoint}
\ContentsText{0pt}{II. Tafeln für die Störungen der Ceres.}
\ContentsText{1.5em}{\textit{Briefwechsel:}}
\ContentsEntry{1.5em}{\textsc{Gauss} an \textsc{Olbers} 1802 December 3}{--- 396}{\ContentsPoint\ContentsPoint\ContentsPoint\ContentsPoint\ContentsPoint\ContentsPoint\ContentsPoint\ContentsPoint\ContentsPoint\ContentsPoint\ContentsPoint\ContentsPoint\ContentsPoint\ContentsPoint\ContentsPoint\ContentsPoint\ContentsPoint\ContentsPoint\ContentsPoint\ContentsPoint}
\ContentsEntry{1.5em}{\textsc{Gauss} an \textsc{Olbers} 1802 December 21}{--- 396}{\ContentsPoint\ContentsPoint\ContentsPoint\ContentsPoint\ContentsPoint\ContentsPoint\ContentsPoint\ContentsPoint\ContentsPoint\ContentsPoint\ContentsPoint\ContentsPoint\ContentsPoint\ContentsPoint\ContentsPoint\ContentsPoint\ContentsPoint\ContentsPoint\ContentsPoint}
\ContentsEntry{1.5em}{\textsc{Gauss} an \textsc{Olbers} 1803 Januar 4}{--- 397}{\ContentsPoint\ContentsPoint\ContentsPoint\ContentsPoint\ContentsPoint\ContentsPoint\ContentsPoint\ContentsPoint\ContentsPoint\ContentsPoint\ContentsPoint\ContentsPoint\ContentsPoint\ContentsPoint\ContentsPoint\ContentsPoint\ContentsPoint\ContentsPoint\ContentsPoint\ContentsPoint\ContentsPoint}
\ContentsEntry{1.5em}{\textsc{Gauss} an \textsc{Olbers} 1803 März 1}{--- 399}{\ContentsPoint\ContentsPoint\ContentsPoint\ContentsPoint\ContentsPoint\ContentsPoint\ContentsPoint\ContentsPoint\ContentsPoint\ContentsPoint\ContentsPoint\ContentsPoint\ContentsPoint\ContentsPoint\ContentsPoint\ContentsPoint\ContentsPoint\ContentsPoint\ContentsPoint\ContentsPoint\ContentsPoint\ContentsPoint}
\ContentsEntry{1.5em}{\textsc{Gauss} an \textsc{Olbers} 1803 April 8}{--- 399}{\ContentsPoint\ContentsPoint\ContentsPoint\ContentsPoint\ContentsPoint\ContentsPoint\ContentsPoint\ContentsPoint\ContentsPoint\ContentsPoint\ContentsPoint\ContentsPoint\ContentsPoint\ContentsPoint\ContentsPoint\ContentsPoint\ContentsPoint\ContentsPoint\ContentsPoint\ContentsPoint\ContentsPoint\ContentsPoint}
\ContentsEntry{1.5em}{\textsc{Gauss} an \textsc{Olbers} 1805 Januar 25}{--- 400}{\ContentsPoint\ContentsPoint\ContentsPoint\ContentsPoint\ContentsPoint\ContentsPoint\ContentsPoint\ContentsPoint\ContentsPoint\ContentsPoint\ContentsPoint\ContentsPoint\ContentsPoint\ContentsPoint\ContentsPoint\ContentsPoint\ContentsPoint\ContentsPoint\ContentsPoint\ContentsPoint\ContentsPoint}
\ContentsText{0pt}{III. Zweite Methode.}
\ContentsEntry{1.5em}{\textit{Briefwechsel:} \textsc{Gauss} an \textsc{Olbers} 1805 Mai 10}{--- 401}{\ContentsPoint\ContentsPoint\ContentsPoint\ContentsPoint\ContentsPoint\ContentsPoint\ContentsPoint\ContentsPoint\ContentsPoint\ContentsPoint\ContentsPoint\ContentsPoint\ContentsPoint\ContentsPoint\ContentsPoint\ContentsPoint\ContentsPoint}

\SourceLeaf{653}
\needspace{4\baselineskip}
\noindent\hfill INHALT.\hfill 647\par
\ContentsEntry{1.5em}{\textit{Nachlass:} Ausführung der Rechnung}{Seite 402}{\ContentsPoint\ContentsPoint\ContentsPoint\ContentsPoint\ContentsPoint\ContentsPoint\ContentsPoint\ContentsPoint\ContentsPoint\ContentsPoint\ContentsPoint\ContentsPoint\ContentsPoint\ContentsPoint\ContentsPoint\ContentsPoint\ContentsPoint\ContentsPoint\ContentsPoint\ContentsPoint}
\ContentsEntry{1.5em}{\textit{Briefwechsel:} \textsc{Gauss} an \textsc{Olbers} 1805 Juli 2}{--- 407}{\ContentsPoint\ContentsPoint\ContentsPoint\ContentsPoint\ContentsPoint\ContentsPoint\ContentsPoint\ContentsPoint\ContentsPoint\ContentsPoint\ContentsPoint\ContentsPoint\ContentsPoint\ContentsPoint\ContentsPoint\ContentsPoint\ContentsPoint}
\ContentsEntry{0pt}{Bemerkungen}{--- 410}{\ContentsPoint\ContentsPoint\ContentsPoint\ContentsPoint\ContentsPoint\ContentsPoint\ContentsPoint\ContentsPoint\ContentsPoint\ContentsPoint\ContentsPoint\ContentsPoint\ContentsPoint\ContentsPoint\ContentsPoint\ContentsPoint\ContentsPoint\ContentsPoint\ContentsPoint\ContentsPoint\ContentsPoint\ContentsPoint\ContentsPoint\ContentsPoint\ContentsPoint\ContentsPoint\ContentsPoint\ContentsPoint\ContentsPoint\ContentsPoint}
\ContentsRule
\ContentsHeading{STÖRUNGEN DER PALLAS.}
\ContentsText{0pt}{I. \textit{Briefwechsel:}}
\ContentsEntry{1.5em}{\textsc{Gauss} an \textsc{Olbers} 1806 Januar 3}{--- 413}{\ContentsPoint\ContentsPoint\ContentsPoint\ContentsPoint\ContentsPoint\ContentsPoint\ContentsPoint\ContentsPoint\ContentsPoint\ContentsPoint\ContentsPoint\ContentsPoint\ContentsPoint\ContentsPoint\ContentsPoint\ContentsPoint\ContentsPoint\ContentsPoint\ContentsPoint\ContentsPoint\ContentsPoint}
\ContentsEntry{1.5em}{\textsc{Gauss} an \textsc{Olbers} 1810 August 6}{--- 413}{\ContentsPoint\ContentsPoint\ContentsPoint\ContentsPoint\ContentsPoint\ContentsPoint\ContentsPoint\ContentsPoint\ContentsPoint\ContentsPoint\ContentsPoint\ContentsPoint\ContentsPoint\ContentsPoint\ContentsPoint\ContentsPoint\ContentsPoint\ContentsPoint\ContentsPoint\ContentsPoint\ContentsPoint}
\ContentsEntry{1.5em}{\textsc{Gauss} an \textsc{Olbers} 1810 October 24}{--- 414}{\ContentsPoint\ContentsPoint\ContentsPoint\ContentsPoint\ContentsPoint\ContentsPoint\ContentsPoint\ContentsPoint\ContentsPoint\ContentsPoint\ContentsPoint\ContentsPoint\ContentsPoint\ContentsPoint\ContentsPoint\ContentsPoint\ContentsPoint\ContentsPoint\ContentsPoint\ContentsPoint}
\ContentsEntry{1.5em}{\textsc{Gauss} an \textsc{Olbers} 1810 November 26}{--- 415}{\ContentsPoint\ContentsPoint\ContentsPoint\ContentsPoint\ContentsPoint\ContentsPoint\ContentsPoint\ContentsPoint\ContentsPoint\ContentsPoint\ContentsPoint\ContentsPoint\ContentsPoint\ContentsPoint\ContentsPoint\ContentsPoint\ContentsPoint\ContentsPoint\ContentsPoint}
\ContentsEntry{1.5em}{\textsc{Gauss} an \textsc{Olbers} 1810 November 30}{--- 417}{\ContentsPoint\ContentsPoint\ContentsPoint\ContentsPoint\ContentsPoint\ContentsPoint\ContentsPoint\ContentsPoint\ContentsPoint\ContentsPoint\ContentsPoint\ContentsPoint\ContentsPoint\ContentsPoint\ContentsPoint\ContentsPoint\ContentsPoint\ContentsPoint\ContentsPoint}
\ContentsEntry{1.5em}{\textsc{Gauss} an \textsc{Olbers} 1810 December 13}{--- 418}{\ContentsPoint\ContentsPoint\ContentsPoint\ContentsPoint\ContentsPoint\ContentsPoint\ContentsPoint\ContentsPoint\ContentsPoint\ContentsPoint\ContentsPoint\ContentsPoint\ContentsPoint\ContentsPoint\ContentsPoint\ContentsPoint\ContentsPoint\ContentsPoint\ContentsPoint\ContentsPoint}
\ContentsEntry{1.5em}{\textsc{Olbers} an \textsc{Gauss} 1810 December 19}{--- 419}{\ContentsPoint\ContentsPoint\ContentsPoint\ContentsPoint\ContentsPoint\ContentsPoint\ContentsPoint\ContentsPoint\ContentsPoint\ContentsPoint\ContentsPoint\ContentsPoint\ContentsPoint\ContentsPoint\ContentsPoint\ContentsPoint\ContentsPoint\ContentsPoint\ContentsPoint\ContentsPoint}
\ContentsEntry{1.5em}{\textsc{Gauss} an \textsc{Schumacher} 1811 Januar 6}{--- 419}{\ContentsPoint\ContentsPoint\ContentsPoint\ContentsPoint\ContentsPoint\ContentsPoint\ContentsPoint\ContentsPoint\ContentsPoint\ContentsPoint\ContentsPoint\ContentsPoint\ContentsPoint\ContentsPoint\ContentsPoint\ContentsPoint\ContentsPoint\ContentsPoint\ContentsPoint}
\ContentsEntry{1.5em}{\textsc{Olbers} an \textsc{Gauss} 1811 Januar 26}{--- 419}{\ContentsPoint\ContentsPoint\ContentsPoint\ContentsPoint\ContentsPoint\ContentsPoint\ContentsPoint\ContentsPoint\ContentsPoint\ContentsPoint\ContentsPoint\ContentsPoint\ContentsPoint\ContentsPoint\ContentsPoint\ContentsPoint\ContentsPoint\ContentsPoint\ContentsPoint\ContentsPoint\ContentsPoint}
\ContentsEntry{1.5em}{\textsc{Schumacher} an \textsc{Gauss} 1811 Juni 31}{--- 419}{\ContentsPoint\ContentsPoint\ContentsPoint\ContentsPoint\ContentsPoint\ContentsPoint\ContentsPoint\ContentsPoint\ContentsPoint\ContentsPoint\ContentsPoint\ContentsPoint\ContentsPoint\ContentsPoint\ContentsPoint\ContentsPoint\ContentsPoint\ContentsPoint\ContentsPoint}
\ContentsEntry{1.5em}{\textsc{Gauss} an \textsc{Olbers} 1811 August 12}{--- 420}{\ContentsPoint\ContentsPoint\ContentsPoint\ContentsPoint\ContentsPoint\ContentsPoint\ContentsPoint\ContentsPoint\ContentsPoint\ContentsPoint\ContentsPoint\ContentsPoint\ContentsPoint\ContentsPoint\ContentsPoint\ContentsPoint\ContentsPoint\ContentsPoint\ContentsPoint\ContentsPoint}
\ContentsEntry{1.5em}{\textsc{Olbers} an \textsc{Gauss} 1812 April 5}{--- 421}{\ContentsPoint\ContentsPoint\ContentsPoint\ContentsPoint\ContentsPoint\ContentsPoint\ContentsPoint\ContentsPoint\ContentsPoint\ContentsPoint\ContentsPoint\ContentsPoint\ContentsPoint\ContentsPoint\ContentsPoint\ContentsPoint\ContentsPoint\ContentsPoint\ContentsPoint\ContentsPoint\ContentsPoint}
\ContentsEntry{1.5em}{\textsc{Gauss} an \textsc{Bessel} 1812 Mai 5}{--- 421}{\ContentsPoint\ContentsPoint\ContentsPoint\ContentsPoint\ContentsPoint\ContentsPoint\ContentsPoint\ContentsPoint\ContentsPoint\ContentsPoint\ContentsPoint\ContentsPoint\ContentsPoint\ContentsPoint\ContentsPoint\ContentsPoint\ContentsPoint\ContentsPoint\ContentsPoint\ContentsPoint\ContentsPoint\ContentsPoint}
\ContentsEntry{1.5em}{\textsc{Olbers} an \textsc{Gauss} 1812 Mai 12}{--- 421}{\ContentsPoint\ContentsPoint\ContentsPoint\ContentsPoint\ContentsPoint\ContentsPoint\ContentsPoint\ContentsPoint\ContentsPoint\ContentsPoint\ContentsPoint\ContentsPoint\ContentsPoint\ContentsPoint\ContentsPoint\ContentsPoint\ContentsPoint\ContentsPoint\ContentsPoint\ContentsPoint\ContentsPoint\ContentsPoint}
\ContentsEntry{1.5em}{\textsc{Gauss} an \textsc{Olbers} 1813 April 8}{--- 422}{\ContentsPoint\ContentsPoint\ContentsPoint\ContentsPoint\ContentsPoint\ContentsPoint\ContentsPoint\ContentsPoint\ContentsPoint\ContentsPoint\ContentsPoint\ContentsPoint\ContentsPoint\ContentsPoint\ContentsPoint\ContentsPoint\ContentsPoint\ContentsPoint\ContentsPoint\ContentsPoint\ContentsPoint\ContentsPoint}
\ContentsEntry{1.5em}{\textsc{Gauss} an \textsc{Olbers} 1813 Juli 2}{--- 422}{\ContentsPoint\ContentsPoint\ContentsPoint\ContentsPoint\ContentsPoint\ContentsPoint\ContentsPoint\ContentsPoint\ContentsPoint\ContentsPoint\ContentsPoint\ContentsPoint\ContentsPoint\ContentsPoint\ContentsPoint\ContentsPoint\ContentsPoint\ContentsPoint\ContentsPoint\ContentsPoint\ContentsPoint\ContentsPoint}
\ContentsEntry{1.5em}{\textsc{Gauss} an \textsc{Olbers} 1814 April 23}{--- 422}{\ContentsPoint\ContentsPoint\ContentsPoint\ContentsPoint\ContentsPoint\ContentsPoint\ContentsPoint\ContentsPoint\ContentsPoint\ContentsPoint\ContentsPoint\ContentsPoint\ContentsPoint\ContentsPoint\ContentsPoint\ContentsPoint\ContentsPoint\ContentsPoint\ContentsPoint\ContentsPoint\ContentsPoint}
\ContentsEntry{1.5em}{\textsc{Gauss} an \textsc{Bessel} 1814 Mai 18}{--- 424}{\ContentsPoint\ContentsPoint\ContentsPoint\ContentsPoint\ContentsPoint\ContentsPoint\ContentsPoint\ContentsPoint\ContentsPoint\ContentsPoint\ContentsPoint\ContentsPoint\ContentsPoint\ContentsPoint\ContentsPoint\ContentsPoint\ContentsPoint\ContentsPoint\ContentsPoint\ContentsPoint\ContentsPoint\ContentsPoint}
\ContentsEntry{1.5em}{\textsc{Gauss} an \textsc{Olbers} 1814 Juni 15}{--- 424}{\ContentsPoint\ContentsPoint\ContentsPoint\ContentsPoint\ContentsPoint\ContentsPoint\ContentsPoint\ContentsPoint\ContentsPoint\ContentsPoint\ContentsPoint\ContentsPoint\ContentsPoint\ContentsPoint\ContentsPoint\ContentsPoint\ContentsPoint\ContentsPoint\ContentsPoint\ContentsPoint\ContentsPoint}
\ContentsEntry{1.5em}{\textsc{Gauss} an \textsc{Olbers} 1814 September 25}{--- 425}{\ContentsPoint\ContentsPoint\ContentsPoint\ContentsPoint\ContentsPoint\ContentsPoint\ContentsPoint\ContentsPoint\ContentsPoint\ContentsPoint\ContentsPoint\ContentsPoint\ContentsPoint\ContentsPoint\ContentsPoint\ContentsPoint\ContentsPoint\ContentsPoint\ContentsPoint}
\ContentsEntry{1.5em}{\textsc{Gauss} an \textsc{Olbers} 1814 December 31}{--- 425}{\ContentsPoint\ContentsPoint\ContentsPoint\ContentsPoint\ContentsPoint\ContentsPoint\ContentsPoint\ContentsPoint\ContentsPoint\ContentsPoint\ContentsPoint\ContentsPoint\ContentsPoint\ContentsPoint\ContentsPoint\ContentsPoint\ContentsPoint\ContentsPoint\ContentsPoint}
\ContentsEntry{1.5em}{\textsc{Olbers} an \textsc{Gauss} 1815 Januar 25}{--- 426}{\ContentsPoint\ContentsPoint\ContentsPoint\ContentsPoint\ContentsPoint\ContentsPoint\ContentsPoint\ContentsPoint\ContentsPoint\ContentsPoint\ContentsPoint\ContentsPoint\ContentsPoint\ContentsPoint\ContentsPoint\ContentsPoint\ContentsPoint\ContentsPoint\ContentsPoint\ContentsPoint\ContentsPoint}
\ContentsEntry{1.5em}{\textsc{Gauss} an \textsc{Olbers} 1816 Januar 8}{--- 426}{\ContentsPoint\ContentsPoint\ContentsPoint\ContentsPoint\ContentsPoint\ContentsPoint\ContentsPoint\ContentsPoint\ContentsPoint\ContentsPoint\ContentsPoint\ContentsPoint\ContentsPoint\ContentsPoint\ContentsPoint\ContentsPoint\ContentsPoint\ContentsPoint\ContentsPoint\ContentsPoint\ContentsPoint}
\ContentsEntry{1.5em}{\textsc{Gauss} an \textsc{Olbers} 1816 Februar 16}{--- 426}{\ContentsPoint\ContentsPoint\ContentsPoint\ContentsPoint\ContentsPoint\ContentsPoint\ContentsPoint\ContentsPoint\ContentsPoint\ContentsPoint\ContentsPoint\ContentsPoint\ContentsPoint\ContentsPoint\ContentsPoint\ContentsPoint\ContentsPoint\ContentsPoint\ContentsPoint\ContentsPoint}
\ContentsEntry{1.5em}{\textsc{Olbers} an \textsc{Gauss} 1816 März 7}{--- 427}{\ContentsPoint\ContentsPoint\ContentsPoint\ContentsPoint\ContentsPoint\ContentsPoint\ContentsPoint\ContentsPoint\ContentsPoint\ContentsPoint\ContentsPoint\ContentsPoint\ContentsPoint\ContentsPoint\ContentsPoint\ContentsPoint\ContentsPoint\ContentsPoint\ContentsPoint\ContentsPoint\ContentsPoint\ContentsPoint}
\ContentsEntry{1.5em}{\textsc{Gauss} an \textsc{Olbers} 1816 Juli 24}{--- 428}{\ContentsPoint\ContentsPoint\ContentsPoint\ContentsPoint\ContentsPoint\ContentsPoint\ContentsPoint\ContentsPoint\ContentsPoint\ContentsPoint\ContentsPoint\ContentsPoint\ContentsPoint\ContentsPoint\ContentsPoint\ContentsPoint\ContentsPoint\ContentsPoint\ContentsPoint\ContentsPoint\ContentsPoint\ContentsPoint}
\ContentsEntry{1.5em}{\textsc{Gauss} an \textsc{Olbers} 1817 Februar 15}{--- 428}{\ContentsPoint\ContentsPoint\ContentsPoint\ContentsPoint\ContentsPoint\ContentsPoint\ContentsPoint\ContentsPoint\ContentsPoint\ContentsPoint\ContentsPoint\ContentsPoint\ContentsPoint\ContentsPoint\ContentsPoint\ContentsPoint\ContentsPoint\ContentsPoint\ContentsPoint\ContentsPoint}
\ContentsEntry{1.5em}{\textsc{Gauss} an \textsc{Olbers} 1817 December 2}{--- 429}{\ContentsPoint\ContentsPoint\ContentsPoint\ContentsPoint\ContentsPoint\ContentsPoint\ContentsPoint\ContentsPoint\ContentsPoint\ContentsPoint\ContentsPoint\ContentsPoint\ContentsPoint\ContentsPoint\ContentsPoint\ContentsPoint\ContentsPoint\ContentsPoint\ContentsPoint\ContentsPoint}
\ContentsEntry{1.5em}{\textsc{Gauss} an \textsc{Encke} 1818 März 25}{--- 430}{\ContentsPoint\ContentsPoint\ContentsPoint\ContentsPoint\ContentsPoint\ContentsPoint\ContentsPoint\ContentsPoint\ContentsPoint\ContentsPoint\ContentsPoint\ContentsPoint\ContentsPoint\ContentsPoint\ContentsPoint\ContentsPoint\ContentsPoint\ContentsPoint\ContentsPoint\ContentsPoint\ContentsPoint}
\ContentsEntry{1.5em}{\textsc{Gauss} an \textsc{Olbers} 1818 März 31}{--- 431}{\ContentsPoint\ContentsPoint\ContentsPoint\ContentsPoint\ContentsPoint\ContentsPoint\ContentsPoint\ContentsPoint\ContentsPoint\ContentsPoint\ContentsPoint\ContentsPoint\ContentsPoint\ContentsPoint\ContentsPoint\ContentsPoint\ContentsPoint\ContentsPoint\ContentsPoint\ContentsPoint\ContentsPoint}
\ContentsEntry{1.5em}{\textsc{Gauss} an \textsc{Schumacher} 1822 März 17}{--- 431}{\ContentsPoint\ContentsPoint\ContentsPoint\ContentsPoint\ContentsPoint\ContentsPoint\ContentsPoint\ContentsPoint\ContentsPoint\ContentsPoint\ContentsPoint\ContentsPoint\ContentsPoint\ContentsPoint\ContentsPoint\ContentsPoint\ContentsPoint\ContentsPoint\ContentsPoint}
\ContentsEntry{1.5em}{\textsc{Gauss} an \textsc{Gerling} 1834 Februar 8}{--- 432}{\ContentsPoint\ContentsPoint\ContentsPoint\ContentsPoint\ContentsPoint\ContentsPoint\ContentsPoint\ContentsPoint\ContentsPoint\ContentsPoint\ContentsPoint\ContentsPoint\ContentsPoint\ContentsPoint\ContentsPoint\ContentsPoint\ContentsPoint\ContentsPoint\ContentsPoint\ContentsPoint}
\ContentsEntry{1.5em}{\textsc{Encke} an \textsc{Gauss} 1834 October 4}{--- 432}{\ContentsPoint\ContentsPoint\ContentsPoint\ContentsPoint\ContentsPoint\ContentsPoint\ContentsPoint\ContentsPoint\ContentsPoint\ContentsPoint\ContentsPoint\ContentsPoint\ContentsPoint\ContentsPoint\ContentsPoint\ContentsPoint\ContentsPoint\ContentsPoint\ContentsPoint\ContentsPoint\ContentsPoint}
\ContentsEntry{1.5em}{\textsc{Gauss} an \textsc{Encke} 1834 October 13}{--- 433}{\ContentsPoint\ContentsPoint\ContentsPoint\ContentsPoint\ContentsPoint\ContentsPoint\ContentsPoint\ContentsPoint\ContentsPoint\ContentsPoint\ContentsPoint\ContentsPoint\ContentsPoint\ContentsPoint\ContentsPoint\ContentsPoint\ContentsPoint\ContentsPoint\ContentsPoint\ContentsPoint}
\ContentsEntry{1.5em}{\textsc{Hansen} an \textsc{Gauss} 1843 Februar 7}{--- 433}{\ContentsPoint\ContentsPoint\ContentsPoint\ContentsPoint\ContentsPoint\ContentsPoint\ContentsPoint\ContentsPoint\ContentsPoint\ContentsPoint\ContentsPoint\ContentsPoint\ContentsPoint\ContentsPoint\ContentsPoint\ContentsPoint\ContentsPoint\ContentsPoint\ContentsPoint\ContentsPoint}
\ContentsEntry{1.5em}{\textsc{Gauss} an \textsc{Hansen} 1843 März 11}{--- 436}{\ContentsPoint\ContentsPoint\ContentsPoint\ContentsPoint\ContentsPoint\ContentsPoint\ContentsPoint\ContentsPoint\ContentsPoint\ContentsPoint\ContentsPoint\ContentsPoint\ContentsPoint\ContentsPoint\ContentsPoint\ContentsPoint\ContentsPoint\ContentsPoint\ContentsPoint\ContentsPoint\ContentsPoint}
\ContentsEntry{1.5em}{\textsc{Gauss} an \textsc{Bessel} 1843 März 21}{--- 437}{\ContentsPoint\ContentsPoint\ContentsPoint\ContentsPoint\ContentsPoint\ContentsPoint\ContentsPoint\ContentsPoint\ContentsPoint\ContentsPoint\ContentsPoint\ContentsPoint\ContentsPoint\ContentsPoint\ContentsPoint\ContentsPoint\ContentsPoint\ContentsPoint\ContentsPoint\ContentsPoint\ContentsPoint}

\SourceLeaf{654}
\needspace{4\baselineskip}
\noindent 648\hfill INHALT.\hfill\mbox{}\par
\ContentsText{0pt}{II. \textit{Nachlass.} Exposition d’une nouvelle méthode de calculer les perturbations planétaires avec l’application au calcul numérique des perturbations du mouvement de Pallas.}
\ContentsEntry{1.5em}{Préface}{Seite 439}{\ContentsPoint\ContentsPoint\ContentsPoint\ContentsPoint\ContentsPoint\ContentsPoint\ContentsPoint\ContentsPoint\ContentsPoint\ContentsPoint\ContentsPoint\ContentsPoint\ContentsPoint\ContentsPoint\ContentsPoint\ContentsPoint\ContentsPoint\ContentsPoint\ContentsPoint\ContentsPoint\ContentsPoint\ContentsPoint\ContentsPoint\ContentsPoint\ContentsPoint\ContentsPoint\ContentsPoint\ContentsPoint\ContentsPoint\ContentsPoint}
\ContentsEntry{1.5em}{Section I. Mouvement elliptique}{--- 440}{\ContentsPoint\ContentsPoint\ContentsPoint\ContentsPoint\ContentsPoint\ContentsPoint\ContentsPoint\ContentsPoint\ContentsPoint\ContentsPoint\ContentsPoint\ContentsPoint\ContentsPoint\ContentsPoint\ContentsPoint\ContentsPoint\ContentsPoint\ContentsPoint\ContentsPoint\ContentsPoint\ContentsPoint}
\ContentsEntry{1.5em}{Section II. Variations instantanées des élémens, produites par les perturbations}{--- 449}{\ContentsPoint\ContentsPoint\ContentsPoint\ContentsPoint}
\ContentsEntry{1.5em}{Section III. Méthode de calculer les variations finies des élémens pendant un tems limité}{--- 460}{}
\ContentsEntry{1.5em}{Section IV. Erster Entwurf. Principes de la détermination du mouvement pour un tems illimité}{--- 465}{\ContentsPoint\ContentsPoint\ContentsPoint\ContentsPoint\ContentsPoint\ContentsPoint\ContentsPoint\ContentsPoint\ContentsPoint\ContentsPoint\ContentsPoint\ContentsPoint\ContentsPoint\ContentsPoint\ContentsPoint\ContentsPoint\ContentsPoint\ContentsPoint\ContentsPoint\ContentsPoint\ContentsPoint\ContentsPoint\ContentsPoint\ContentsPoint\ContentsPoint\ContentsPoint\ContentsPoint\ContentsPoint\ContentsPoint}
\ContentsEntry{1.5em}{Section IV. Zweiter Entwurf. Développement des fonctions périodiques en séries}{--- 469}{\ContentsPoint\ContentsPoint\ContentsPoint}
\ContentsText{0pt}{III. \textit{Nachlass.} Specielle Störungen der Pallas durch Jupiter. Erste Rechnung.}
\ContentsEntry{1.5em}{Ableitung der Formeln}{--- 473}{\ContentsPoint\ContentsPoint\ContentsPoint\ContentsPoint\ContentsPoint\ContentsPoint\ContentsPoint\ContentsPoint\ContentsPoint\ContentsPoint\ContentsPoint\ContentsPoint\ContentsPoint\ContentsPoint\ContentsPoint\ContentsPoint\ContentsPoint\ContentsPoint\ContentsPoint\ContentsPoint\ContentsPoint\ContentsPoint\ContentsPoint\ContentsPoint}
\ContentsEntry{1.5em}{Zum Grunde gelegte Elementensysteme}{--- 474}{\ContentsPoint\ContentsPoint\ContentsPoint\ContentsPoint\ContentsPoint\ContentsPoint\ContentsPoint\ContentsPoint\ContentsPoint\ContentsPoint\ContentsPoint\ContentsPoint\ContentsPoint\ContentsPoint\ContentsPoint\ContentsPoint\ContentsPoint\ContentsPoint\ContentsPoint}
\ContentsEntry{1.5em}{Berechnung der Störungen}{--- 475}{\ContentsPoint\ContentsPoint\ContentsPoint\ContentsPoint\ContentsPoint\ContentsPoint\ContentsPoint\ContentsPoint\ContentsPoint\ContentsPoint\ContentsPoint\ContentsPoint\ContentsPoint\ContentsPoint\ContentsPoint\ContentsPoint\ContentsPoint\ContentsPoint\ContentsPoint\ContentsPoint\ContentsPoint\ContentsPoint\ContentsPoint}
\ContentsEntry{1.5em}{Bestimmung der Normalelemente}{--- 479}{\ContentsPoint\ContentsPoint\ContentsPoint\ContentsPoint\ContentsPoint\ContentsPoint\ContentsPoint\ContentsPoint\ContentsPoint\ContentsPoint\ContentsPoint\ContentsPoint\ContentsPoint\ContentsPoint\ContentsPoint\ContentsPoint\ContentsPoint\ContentsPoint\ContentsPoint\ContentsPoint\ContentsPoint}
\ContentsText{0pt}{IV. \textit{Nachlass.} Specielle Störungen der Pallas durch Jupiter. Zweite Rechnung.}
\ContentsEntry{1.5em}{Berechnung der Störungen}{--- 483}{\ContentsPoint\ContentsPoint\ContentsPoint\ContentsPoint\ContentsPoint\ContentsPoint\ContentsPoint\ContentsPoint\ContentsPoint\ContentsPoint\ContentsPoint\ContentsPoint\ContentsPoint\ContentsPoint\ContentsPoint\ContentsPoint\ContentsPoint\ContentsPoint\ContentsPoint\ContentsPoint\ContentsPoint\ContentsPoint\ContentsPoint}
\ContentsEntry{1.5em}{Bestimmung der Normalelemente}{--- 486}{\ContentsPoint\ContentsPoint\ContentsPoint\ContentsPoint\ContentsPoint\ContentsPoint\ContentsPoint\ContentsPoint\ContentsPoint\ContentsPoint\ContentsPoint\ContentsPoint\ContentsPoint\ContentsPoint\ContentsPoint\ContentsPoint\ContentsPoint\ContentsPoint\ContentsPoint\ContentsPoint\ContentsPoint}
\ContentsText{0pt}{V. \textit{Nachlass.} Allgemeine Störungen der Pallas durch Jupiter. Erste Rechnung.}
\ContentsEntry{1.5em}{Formeln für die Variation der Elemente}{--- 489}{\ContentsPoint\ContentsPoint\ContentsPoint\ContentsPoint\ContentsPoint\ContentsPoint\ContentsPoint\ContentsPoint\ContentsPoint\ContentsPoint\ContentsPoint\ContentsPoint\ContentsPoint\ContentsPoint\ContentsPoint\ContentsPoint\ContentsPoint\ContentsPoint}
\ContentsEntry{1.5em}{Entwickelung einer periodischen Function in eine Reihe}{--- 491}{\ContentsPoint\ContentsPoint\ContentsPoint\ContentsPoint\ContentsPoint\ContentsPoint\ContentsPoint\ContentsPoint\ContentsPoint\ContentsPoint\ContentsPoint\ContentsPoint\ContentsPoint}
\ContentsEntry{1.5em}{Praktische Anwendung des Vorigen}{--- 496}{\ContentsPoint\ContentsPoint\ContentsPoint\ContentsPoint\ContentsPoint\ContentsPoint\ContentsPoint\ContentsPoint\ContentsPoint\ContentsPoint\ContentsPoint\ContentsPoint\ContentsPoint\ContentsPoint\ContentsPoint\ContentsPoint\ContentsPoint\ContentsPoint\ContentsPoint\ContentsPoint}
\ContentsEntry{1.5em}{Rechnungsbeispiel}{--- 500}{\ContentsPoint\ContentsPoint\ContentsPoint\ContentsPoint\ContentsPoint\ContentsPoint\ContentsPoint\ContentsPoint\ContentsPoint\ContentsPoint\ContentsPoint\ContentsPoint\ContentsPoint\ContentsPoint\ContentsPoint\ContentsPoint\ContentsPoint\ContentsPoint\ContentsPoint\ContentsPoint\ContentsPoint\ContentsPoint\ContentsPoint\ContentsPoint\ContentsPoint\ContentsPoint}
\ContentsEntry{1.5em}{Entwickelung einer periodischen Function zweier Veränderlicher}{--- 505}{\ContentsPoint\ContentsPoint\ContentsPoint\ContentsPoint\ContentsPoint\ContentsPoint\ContentsPoint\ContentsPoint\ContentsPoint\ContentsPoint}
\ContentsEntry{1.5em}{Berechnung und Entwickelung der störenden Kräfte}{--- 507}{\ContentsPoint\ContentsPoint\ContentsPoint\ContentsPoint\ContentsPoint\ContentsPoint\ContentsPoint\ContentsPoint\ContentsPoint\ContentsPoint\ContentsPoint\ContentsPoint\ContentsPoint\ContentsPoint}
\ContentsEntry{1.5em}{Störungen der Neigung und der Länge des Knoten}{--- 513}{\ContentsPoint\ContentsPoint\ContentsPoint\ContentsPoint\ContentsPoint\ContentsPoint\ContentsPoint\ContentsPoint\ContentsPoint\ContentsPoint\ContentsPoint\ContentsPoint\ContentsPoint\ContentsPoint}
\ContentsEntry{1.5em}{Störungen des Perihels}{--- 519}{\ContentsPoint\ContentsPoint\ContentsPoint\ContentsPoint\ContentsPoint\ContentsPoint\ContentsPoint\ContentsPoint\ContentsPoint\ContentsPoint\ContentsPoint\ContentsPoint\ContentsPoint\ContentsPoint\ContentsPoint\ContentsPoint\ContentsPoint\ContentsPoint\ContentsPoint\ContentsPoint\ContentsPoint\ContentsPoint\ContentsPoint\ContentsPoint\ContentsPoint}
\ContentsEntry{1.5em}{Störungen des Logarithmen des halben Parameters}{--- 521}{\ContentsPoint\ContentsPoint\ContentsPoint\ContentsPoint\ContentsPoint\ContentsPoint\ContentsPoint\ContentsPoint\ContentsPoint\ContentsPoint\ContentsPoint\ContentsPoint\ContentsPoint\ContentsPoint\ContentsPoint}
\ContentsEntry{1.5em}{Störungen des Logarithmen der halben grossen Axe}{--- 522}{\ContentsPoint\ContentsPoint\ContentsPoint\ContentsPoint\ContentsPoint\ContentsPoint\ContentsPoint\ContentsPoint\ContentsPoint\ContentsPoint\ContentsPoint\ContentsPoint\ContentsPoint\ContentsPoint}
\ContentsEntry{1.5em}{Störungen des ersten Theils der Epoche}{--- 524}{\ContentsPoint\ContentsPoint\ContentsPoint\ContentsPoint\ContentsPoint\ContentsPoint\ContentsPoint\ContentsPoint\ContentsPoint\ContentsPoint\ContentsPoint\ContentsPoint\ContentsPoint\ContentsPoint\ContentsPoint\ContentsPoint\ContentsPoint\ContentsPoint}
\ContentsEntry{1.5em}{Störungen des zweiten Theils der Epoche}{--- 525}{\ContentsPoint\ContentsPoint\ContentsPoint\ContentsPoint\ContentsPoint\ContentsPoint\ContentsPoint\ContentsPoint\ContentsPoint\ContentsPoint\ContentsPoint\ContentsPoint\ContentsPoint\ContentsPoint\ContentsPoint\ContentsPoint\ContentsPoint\ContentsPoint}
\ContentsEntry{1.5em}{Bestimmung der Normalelemente}{--- 526}{\ContentsPoint\ContentsPoint\ContentsPoint\ContentsPoint\ContentsPoint\ContentsPoint\ContentsPoint\ContentsPoint\ContentsPoint\ContentsPoint\ContentsPoint\ContentsPoint\ContentsPoint\ContentsPoint\ContentsPoint\ContentsPoint\ContentsPoint\ContentsPoint\ContentsPoint\ContentsPoint\ContentsPoint}
\ContentsText{0pt}{VI. \textit{Nachlass.} Allgemeine Störungen der Pallas durch Jupiter. Zweite Rechnung.}
\ContentsEntry{1.5em}{Berechnung und Entwickelung der störenden Kräfte; Formeln für die Variation der Elemente}{--- 529}{}
\ContentsEntry{1.5em}{Störungen der Neigung und der Länge des Knoten}{--- 542}{\ContentsPoint\ContentsPoint\ContentsPoint\ContentsPoint\ContentsPoint\ContentsPoint\ContentsPoint\ContentsPoint\ContentsPoint\ContentsPoint\ContentsPoint\ContentsPoint\ContentsPoint\ContentsPoint}
\ContentsEntry{1.5em}{Störungen des Perihels}{--- 545}{\ContentsPoint\ContentsPoint\ContentsPoint\ContentsPoint\ContentsPoint\ContentsPoint\ContentsPoint\ContentsPoint\ContentsPoint\ContentsPoint\ContentsPoint\ContentsPoint\ContentsPoint\ContentsPoint\ContentsPoint\ContentsPoint\ContentsPoint\ContentsPoint\ContentsPoint\ContentsPoint\ContentsPoint\ContentsPoint\ContentsPoint\ContentsPoint\ContentsPoint}
\ContentsEntry{1.5em}{Störungen des Excentricitätswinkels}{--- 548}{\ContentsPoint\ContentsPoint\ContentsPoint\ContentsPoint\ContentsPoint\ContentsPoint\ContentsPoint\ContentsPoint\ContentsPoint\ContentsPoint\ContentsPoint\ContentsPoint\ContentsPoint\ContentsPoint\ContentsPoint\ContentsPoint\ContentsPoint\ContentsPoint\ContentsPoint\ContentsPoint}
\ContentsEntry{1.5em}{Störungen des Logarithmen der halben grossen Axe}{--- 549}{\ContentsPoint\ContentsPoint\ContentsPoint\ContentsPoint\ContentsPoint\ContentsPoint\ContentsPoint\ContentsPoint\ContentsPoint\ContentsPoint\ContentsPoint\ContentsPoint\ContentsPoint\ContentsPoint}
\ContentsEntry{1.5em}{Störungen des ersten Theils der Epoche}{--- 550}{\ContentsPoint\ContentsPoint\ContentsPoint\ContentsPoint\ContentsPoint\ContentsPoint\ContentsPoint\ContentsPoint\ContentsPoint\ContentsPoint\ContentsPoint\ContentsPoint\ContentsPoint\ContentsPoint\ContentsPoint\ContentsPoint\ContentsPoint\ContentsPoint}
\ContentsEntry{1.5em}{Störungen des zweiten Theils der Epoche}{--- 552}{\ContentsPoint\ContentsPoint\ContentsPoint\ContentsPoint\ContentsPoint\ContentsPoint\ContentsPoint\ContentsPoint\ContentsPoint\ContentsPoint\ContentsPoint\ContentsPoint\ContentsPoint\ContentsPoint\ContentsPoint\ContentsPoint\ContentsPoint\ContentsPoint}
\ContentsEntry{1.5em}{Vereinigung der Störungen der Epoche}{--- 554}{\ContentsPoint\ContentsPoint\ContentsPoint\ContentsPoint\ContentsPoint\ContentsPoint\ContentsPoint\ContentsPoint\ContentsPoint\ContentsPoint\ContentsPoint\ContentsPoint\ContentsPoint\ContentsPoint\ContentsPoint\ContentsPoint\ContentsPoint\ContentsPoint\ContentsPoint}
\ContentsEntry{1.5em}{Praktische Integration der vom Argument \(18\ContentsJupiter-7\ContentsPallas\) abhängenden Gleichungen}{--- 557}{\ContentsPoint\ContentsPoint\ContentsPoint\ContentsPoint\ContentsPoint}
\ContentsEntry{1.5em}{Rationales Verhältniss der Umlaufszeiten von Pallas und Jupiter}{--- 557}{\ContentsPoint\ContentsPoint\ContentsPoint\ContentsPoint\ContentsPoint\ContentsPoint\ContentsPoint\ContentsPoint\ContentsPoint\ContentsPoint}

\SourceLeaf{655}
\needspace{4\baselineskip}
\noindent\hfill INHALT.\hfill 649\par
\ContentsEntry{1.5em}{Bestimmung der Normalelemente}{Seite 559}{\ContentsPoint\ContentsPoint\ContentsPoint\ContentsPoint\ContentsPoint\ContentsPoint\ContentsPoint\ContentsPoint\ContentsPoint\ContentsPoint\ContentsPoint\ContentsPoint\ContentsPoint\ContentsPoint\ContentsPoint\ContentsPoint\ContentsPoint\ContentsPoint\ContentsPoint\ContentsPoint\ContentsPoint}
\ContentsEntry{1.5em}{Verbesserung der Jupitermasse}{--- 561}{\ContentsPoint\ContentsPoint\ContentsPoint\ContentsPoint\ContentsPoint\ContentsPoint\ContentsPoint\ContentsPoint\ContentsPoint\ContentsPoint\ContentsPoint\ContentsPoint\ContentsPoint\ContentsPoint\ContentsPoint\ContentsPoint\ContentsPoint\ContentsPoint\ContentsPoint\ContentsPoint\ContentsPoint\ContentsPoint}
\ContentsText{0pt}{VII. Tafeln für die Störungen der Pallas durch Jupiter.}
\ContentsEntry{1.5em}{\textit{Briefwechsel:} Proberechnung: \textsc{Gauss} an \textsc{Encke} 1816 August 5}{--- 565}{\ContentsPoint\ContentsPoint\ContentsPoint\ContentsPoint\ContentsPoint\ContentsPoint\ContentsPoint\ContentsPoint\ContentsPoint\ContentsPoint}
\ContentsEntry{1.5em}{Formeln zur Berechnung der Tafeln}{--- 567}{\ContentsPoint\ContentsPoint\ContentsPoint\ContentsPoint\ContentsPoint\ContentsPoint\ContentsPoint\ContentsPoint\ContentsPoint\ContentsPoint\ContentsPoint\ContentsPoint\ContentsPoint\ContentsPoint\ContentsPoint\ContentsPoint\ContentsPoint\ContentsPoint\ContentsPoint\ContentsPoint}
\ContentsText{1.5em}{\textit{Briefwechsel:} Schema für die Ausführung der Rechnungen:}
\ContentsEntry{7em}{\textsc{Gauss} an \textsc{Encke} 1816 November 7 und 1877 März 3}{--- 568}{\ContentsPoint\ContentsPoint\ContentsPoint\ContentsPoint\ContentsPoint\ContentsPoint}
\ContentsEntry{1.5em}{\textit{Nachlass:} Tafeln}{--- 572}{\ContentsPoint\ContentsPoint\ContentsPoint\ContentsPoint\ContentsPoint\ContentsPoint\ContentsPoint\ContentsPoint\ContentsPoint\ContentsPoint\ContentsPoint\ContentsPoint\ContentsPoint\ContentsPoint\ContentsPoint\ContentsPoint\ContentsPoint\ContentsPoint\ContentsPoint\ContentsPoint\ContentsPoint\ContentsPoint\ContentsPoint\ContentsPoint\ContentsPoint\ContentsPoint\ContentsPoint}
\ContentsText{0pt}{VIII. Störungen der Pallas durch Saturn, berechnet von \textsc{Nicolai}.}
\ContentsText{1.5em}{\textit{Nachlass.}}
\ContentsEntry{1.5em}{Formeln für die Variation der Elemente}{--- 578}{\ContentsPoint\ContentsPoint\ContentsPoint\ContentsPoint\ContentsPoint\ContentsPoint\ContentsPoint\ContentsPoint\ContentsPoint\ContentsPoint\ContentsPoint\ContentsPoint\ContentsPoint\ContentsPoint\ContentsPoint\ContentsPoint\ContentsPoint\ContentsPoint\ContentsPoint}
\ContentsText{1.5em}{\textit{Briefwechsel.}}
\ContentsText{1.5em}{Störungen des Logarithmen der halben grossen Axe, des ersten Theils der Epoche, der Neigung und der Länge des Knoten:}
\ContentsEntry{10em}{\textsc{Nicolai} an \textsc{Gauss} 1814 März 19}{--- 579}{\ContentsPoint\ContentsPoint\ContentsPoint\ContentsPoint\ContentsPoint\ContentsPoint\ContentsPoint\ContentsPoint\ContentsPoint}
\ContentsText{1.5em}{Störungen des Excentricitätswinkels und der Länge des Perihels:}
\ContentsEntry{10em}{\textsc{Nicolai} an \textsc{Gauss} 1814 am Charfreitage}{--- 581}{\ContentsPoint\ContentsPoint\ContentsPoint\ContentsPoint\ContentsPoint\ContentsPoint}
\ContentsEntry{1.5em}{Störungen des zweiten Theils der Epoche: \textsc{Nicolai} an \textsc{Gauss} 1814 April 22}{--- 583}{\ContentsPoint\ContentsPoint\ContentsPoint\ContentsPoint\ContentsPoint}
\ContentsEntry{1.5em}{Controlle der Rechnungen: \textsc{Nicolai} an \textsc{Gauss} 1814 August 6}{--- 584}{\ContentsPoint\ContentsPoint\ContentsPoint\ContentsPoint\ContentsPoint\ContentsPoint\ContentsPoint\ContentsPoint\ContentsPoint\ContentsPoint\ContentsPoint}
\ContentsEntry{1.5em}{Vereinigung beider Theile der Epoche: \textsc{Nicolai} an \textsc{Gauss} 1815 März 18 und Juli 17}{--- 585}{\ContentsPoint\ContentsPoint}
\ContentsText{0pt}{IX. \textit{Nachlass.} Störungen der Pallas durch Mars.}
\ContentsEntry{1.5em}{Formeln für die Variation der Elemente; Berechnung der störenden Kräfte}{--- 587}{\ContentsPoint\ContentsPoint\ContentsPoint\ContentsPoint\ContentsPoint\ContentsPoint}
\ContentsEntry{1.5em}{Entwickelung der störenden Kräfte}{--- 591}{\ContentsPoint\ContentsPoint\ContentsPoint\ContentsPoint\ContentsPoint\ContentsPoint\ContentsPoint\ContentsPoint\ContentsPoint\ContentsPoint\ContentsPoint\ContentsPoint\ContentsPoint\ContentsPoint\ContentsPoint\ContentsPoint\ContentsPoint\ContentsPoint\ContentsPoint\ContentsPoint}
\ContentsEntry{1.5em}{Störungen des Logarithmen der halben grossen Axe, des ersten Theils der Epoche und des Excentricitätswinkels}{--- 593}{\ContentsPoint\ContentsPoint\ContentsPoint\ContentsPoint\ContentsPoint\ContentsPoint\ContentsPoint\ContentsPoint\ContentsPoint\ContentsPoint\ContentsPoint\ContentsPoint\ContentsPoint\ContentsPoint\ContentsPoint\ContentsPoint\ContentsPoint\ContentsPoint\ContentsPoint\ContentsPoint\ContentsPoint\ContentsPoint\ContentsPoint\ContentsPoint}
\ContentsEntry{1.5em}{Über die Convergenz der Reihen für die störenden Kräfte; andere Methode ihrer Entwickelung durch Einführung eines Factors \(\lambda\)}{--- 594}{\ContentsPoint\ContentsPoint\ContentsPoint\ContentsPoint\ContentsPoint\ContentsPoint\ContentsPoint\ContentsPoint\ContentsPoint\ContentsPoint\ContentsPoint\ContentsPoint\ContentsPoint\ContentsPoint\ContentsPoint\ContentsPoint\ContentsPoint}
\ContentsEntry{1.5em}{Bestimmung des Factors \(\lambda\)}{--- 594}{\ContentsPoint\ContentsPoint\ContentsPoint\ContentsPoint\ContentsPoint\ContentsPoint\ContentsPoint\ContentsPoint\ContentsPoint\ContentsPoint\ContentsPoint\ContentsPoint\ContentsPoint\ContentsPoint\ContentsPoint\ContentsPoint\ContentsPoint\ContentsPoint\ContentsPoint\ContentsPoint\ContentsPoint\ContentsPoint\ContentsPoint}
\ContentsEntry{1.5em}{Aufstellung der Reihen für die störenden Kräfte}{--- 596}{\ContentsPoint\ContentsPoint\ContentsPoint\ContentsPoint\ContentsPoint\ContentsPoint\ContentsPoint\ContentsPoint\ContentsPoint\ContentsPoint\ContentsPoint\ContentsPoint\ContentsPoint\ContentsPoint\ContentsPoint}
\ContentsEntry{1.5em}{Andere Methode zur Bestimmung des Factors \(\lambda\); Auflösung der Gleichung \(a+b\cos2u=\cos(u+D)\)}{--- 599}{\ContentsPoint\ContentsPoint\ContentsPoint\ContentsPoint\ContentsPoint\ContentsPoint\ContentsPoint\ContentsPoint\ContentsPoint\ContentsPoint\ContentsPoint\ContentsPoint\ContentsPoint\ContentsPoint\ContentsPoint\ContentsPoint\ContentsPoint\ContentsPoint\ContentsPoint\ContentsPoint\ContentsPoint\ContentsPoint\ContentsPoint\ContentsPoint\ContentsPoint\ContentsPoint\ContentsPoint}
\ContentsEntry{0pt}{Bemerkungen}{--- 601}{\ContentsPoint\ContentsPoint\ContentsPoint\ContentsPoint\ContentsPoint\ContentsPoint\ContentsPoint\ContentsPoint\ContentsPoint\ContentsPoint\ContentsPoint\ContentsPoint\ContentsPoint\ContentsPoint\ContentsPoint\ContentsPoint\ContentsPoint\ContentsPoint\ContentsPoint\ContentsPoint\ContentsPoint\ContentsPoint\ContentsPoint\ContentsPoint\ContentsPoint\ContentsPoint\ContentsPoint\ContentsPoint\ContentsPoint\ContentsPoint}
\ContentsRule
\ContentsHeading{THEORIE DER BEWEGUNG DES MONDES.}
\ContentsText{0pt}{\textit{Nachlass.}}
\ContentsEntry{1.5em}{Erster Abschnitt. Fundamentalgleichungen für die Bewegung des Mondes}{--- 613}{\ContentsPoint\ContentsPoint\ContentsPoint\ContentsPoint\ContentsPoint\ContentsPoint}
\ContentsEntry{1.5em}{Zweiter Abschnitt. Integration der Fundamentalgleichungen mit Beiseitesetzung der Störungskräfte der Sonne}{--- 617}{\ContentsPoint\ContentsPoint\ContentsPoint\ContentsPoint\ContentsPoint\ContentsPoint\ContentsPoint\ContentsPoint\ContentsPoint\ContentsPoint\ContentsPoint\ContentsPoint\ContentsPoint\ContentsPoint\ContentsPoint\ContentsPoint\ContentsPoint\ContentsPoint\ContentsPoint\ContentsPoint\ContentsPoint\ContentsPoint\ContentsPoint\ContentsPoint\ContentsPoint}
\medskip\noindent VII.\hfill 82\par

\SourceLeaf{656}
\needspace{4\baselineskip}
\noindent 650\hfill INHALT.\hfill\mbox{}\par
\ContentsEntry{1.5em}{Dritter Abschnitt. Erste Annäherung}{Seite 620}{\ContentsPoint\ContentsPoint\ContentsPoint\ContentsPoint\ContentsPoint\ContentsPoint\ContentsPoint\ContentsPoint\ContentsPoint\ContentsPoint\ContentsPoint\ContentsPoint\ContentsPoint\ContentsPoint\ContentsPoint\ContentsPoint\ContentsPoint\ContentsPoint\ContentsPoint}
\ContentsEntry{1.5em}{Vierter Abschnitt. Zweite Annäherung zur Breite}{--- 626}{\ContentsPoint\ContentsPoint\ContentsPoint\ContentsPoint\ContentsPoint\ContentsPoint\ContentsPoint\ContentsPoint\ContentsPoint\ContentsPoint\ContentsPoint\ContentsPoint\ContentsPoint\ContentsPoint\ContentsPoint}
\ContentsEntry{1.5em}{Fünfter Abschnitt. Zweite Annäherung zur Länge}{--- 634}{\ContentsPoint\ContentsPoint\ContentsPoint\ContentsPoint\ContentsPoint\ContentsPoint\ContentsPoint\ContentsPoint\ContentsPoint\ContentsPoint\ContentsPoint\ContentsPoint\ContentsPoint\ContentsPoint\ContentsPoint}
\ContentsEntry{0pt}{Bemerkungen}{--- 640}{\ContentsPoint\ContentsPoint\ContentsPoint\ContentsPoint\ContentsPoint\ContentsPoint\ContentsPoint\ContentsPoint\ContentsPoint\ContentsPoint\ContentsPoint\ContentsPoint\ContentsPoint\ContentsPoint\ContentsPoint\ContentsPoint\ContentsPoint\ContentsPoint\ContentsPoint\ContentsPoint\ContentsPoint\ContentsPoint\ContentsPoint\ContentsPoint\ContentsPoint\ContentsPoint\ContentsPoint\ContentsPoint\ContentsPoint\ContentsPoint}
\ContentsRule
\ContentsEntry{0pt}{Berichtigungen und Zusätze}{--- 642}{\ContentsPoint\ContentsPoint\ContentsPoint\ContentsPoint\ContentsPoint\ContentsPoint\ContentsPoint\ContentsPoint\ContentsPoint\ContentsPoint\ContentsPoint\ContentsPoint\ContentsPoint\ContentsPoint\ContentsPoint\ContentsPoint\ContentsPoint\ContentsPoint\ContentsPoint\ContentsPoint\ContentsPoint\ContentsPoint\ContentsPoint\ContentsPoint\ContentsPoint}
\ContentsEntry{0pt}{Bemerkungen zum siebenten Bande}{--- 643}{\ContentsPoint\ContentsPoint\ContentsPoint\ContentsPoint\ContentsPoint\ContentsPoint\ContentsPoint\ContentsPoint\ContentsPoint\ContentsPoint\ContentsPoint\ContentsPoint\ContentsPoint\ContentsPoint\ContentsPoint\ContentsPoint\ContentsPoint\ContentsPoint\ContentsPoint\ContentsPoint\ContentsPoint\ContentsPoint}
\begin{center}\rule{.28\linewidth}{.4pt}\end{center}
\bigskip
\begin{center}\small Göttingen, Druck der Dieterichschen Univ.-Buchdruckerei (W. Fr. Kaestner).\end{center}
\endgroup

\clearpage
\endgroup
